\section{Threshold and Rate-Complexity Corollaries}\label{sec:rate-corollaries}
%==============================================================================


The main graph-capacity arc characterizes the structure of exact failure under partial views. This section records one downstream operational consequence at the opposite boundary: independent rate governs manual synchronization cost in the same finite deterministic model.

\subsection{Cost Model}\label{sec:cost-model}

\begin{definition}[Modification Cost Model]\label{def:cost-model}
Let $\delta_F$ be a modification to fact $F$ in encoding system $C$. The \emph{effective modification complexity} $M_{\mathrm{effective}}(C,\delta_F)$ is the number of locations that must be edited manually to preserve correctness after applying $\delta_F$.
\end{definition}

Only manual edits are counted. Derived locations updated automatically by the host system contribute zero effective cost.

\subsection{Upper Bound at Unit Independent Rate}\label{sec:upper-bound}

\begin{theorem}[Rate-1 Upper Bound]\label{thm:upper-bound}
If $\text{DOF}(C,F)=1$, then
\[
M_{\mathrm{effective}}(C,\delta_F)=O(1).
\]
\end{theorem}
\leanmeta{\LH{BND1}}

\begin{proof}
When $\text{DOF}(C,F)=1$, there is exactly one independently writable source location for $F$. Updating that location is one manual edit; every other representation of $F$ is derived and is updated automatically. The number of manual edits therefore stays bounded by a constant independent of the number of visible encodings.
\end{proof}

\subsection{Lower Bound Above the Threshold}\label{sec:lower-bound}

\begin{theorem}[Above-Threshold Lower Bound]\label{thm:lower-bound}
If fact $F$ is encoded at $n$ independent locations, then
\[
M_{\mathrm{effective}}(C,\delta_F)=\Omega(n).
\]
\end{theorem}
\leanmeta{\LH{BND2}}

\begin{proof}
Each independent location can retain an outdated value unless it is edited directly. After a source update, every one of the $n$ independent locations can therefore witness a distinct stale copy of the same fact. Any exact repair strategy must bring each such location back into agreement, and in the worst case no single manual action can repair two independently writable locations at once. Hence the effective modification complexity grows at least linearly in $n$.
\end{proof}

\begin{lemma}[Information-Constrained DOF Lower Bound]\label{lem:info-dof}
If the available side-information mechanism exposes at most $2^L$ distinguishable tags while the latent ambiguity class has size $K>2^L$, then no zero-error resolver can identify the true state from those tags alone. Any architecture that still guarantees exact correctness must therefore rely on additional independently accessible discriminating support, moving the system above unit rate.
\end{lemma}
\leanmeta{\LH{CIA4}}

\begin{proof}
Theorem~\ref{thm:fano-converse} rules out exact recovery from the available tag budget. If exact correctness is nevertheless required, the architecture must supplement the insufficient side-information channel with additional independently accessible information. In the present model, that means leaving the rate-$1$ regime and paying the synchronization cost of Theorem~\ref{thm:lower-bound}.
\end{proof}

\subsection{The Unbounded Gap}\label{sec:gap}

\begin{theorem}[Unbounded Gap]\label{thm:unbounded-gap}
The ratio between the modification cost of an architecture with $n$ independent encodings and the modification cost of a rate-$1$ architecture diverges:
\[
\lim_{n\to\infty}\frac{M_{\mathrm{DOF}>1}(n)}{M_{\mathrm{DOF}=1}}=\infty.
\]
\end{theorem}
\leanmeta{\LH{BND3}}

\begin{proof}
By Theorem~\ref{thm:upper-bound}, a rate-$1$ architecture has constant effective modification complexity. By Theorem~\ref{thm:lower-bound}, an architecture with $n$ independent encodings incurs linear cost in $n$. The ratio therefore grows without bound.
\end{proof}

Once a fact is duplicated across many independently writable locations, the cost of preserving exact consistency scales with the number of independent copies, not merely with the visibility of the fact.
